\documentclass[11pt,a4paper]{article}

% -------------------- Langue / encodage --------------------
\usepackage[T1]{fontenc}
\usepackage[french]{babel}
\usepackage{array}

% -------------------- Mise en page --------------------
\usepackage{geometry}
\geometry{
  a4paper,
  landscape,
  left=1.6cm,
  right=1.6cm,
  top=1.0cm,
  bottom=1.2cm
}

% -------------------- Maths / mise en forme --------------------
\usepackage{amsmath,amssymb}
\usepackage{parskip}
\usepackage[explicit]{titlesec}
\usepackage[normalem]{ulem}
\usepackage{titlesec}

\usepackage{fancyhdr}
\usepackage{lastpage}
\usepackage[hidelinks]{hyperref}
\usepackage{enumitem}
\usepackage[table]{xcolor}
\usepackage{paracol}

% -------------------- TikZ --------------------
\usepackage{tikz}
\usetikzlibrary{arrows.meta,calc,intersections}

% -------------------- Variables --------------------
\newcommand{\FicheTitre}{Mécanique dans un champ de gravitation -- Fiche de cours}
\newcommand{\Niveau}{Terminale générale -- Physique}
\newcommand{\Annee}{Année scolaire 2025/2026}

\newcommand{\SiteURL}{https://naoufel-coursparticuliers.netlify.app/}
\newcommand{\SiteTexte}{naoufel-coursparticuliers.netlify.app}

% -------------------- Footer --------------------
\pagestyle{fancy}
\fancyhf{}
\lfoot{\textit{\Niveau \;--\; \Annee}}
\cfoot{\thepage/\pageref{LastPage}}
\rfoot{\href{\SiteURL}{\texttt{\SiteTexte}}}
\renewcommand{\headrulewidth}{0pt}
\renewcommand{\footrulewidth}{0.4pt}

% -------------------- Spacing --------------------
\setlength{\parskip}{4pt}
\setlength{\parindent}{0pt}
\setlist[itemize]{leftmargin=*, itemsep=2pt, topsep=2pt}

% -------------------- Titres --------------------
\titleformat{\section}[hang]
  {\normalsize\bfseries}
  {\textcolor{red}{\thesection.}}
  {0.2em}
  {\textcolor{red}{\uline{#1}}}

\titleformat{\subsection}[hang]
  {\normalsize\bfseries}
  {\hspace*{1.2em}\textcolor{green!60!black}{\thesubsection}}
  {0.4em}
  {\textcolor{green!60!black}{\uline{#1}}}

% -------------------- Bandeau titre --------------------
\newcommand{\TitreBandeau}[1]{%
  \noindent\hrule
  \vspace{0.55em}
  \begin{center}
    {\LARGE\bfseries #1}
  \end{center}
  \vspace{0.55em}
  \noindent\hrule
  \vspace{0.9em}
}

% -------------------- Helpers --------------------
\newcommand{\defi}{\textbf{Définition :} }
\newcommand{\prop}{\textbf{Propriété :} }
\newcommand{\rem}{\textbf{Remarque :} }


% ============================
% Schéma : champ de gravitation
% ============================
\newcommand{\SchemaGravitation}{%
\begin{center}
\resizebox{0.5\linewidth}{!}{%
\begin{tikzpicture}[>=Latex, every node/.style={font=\small}]

% Astre
\fill[gray!20] (0,0) circle (0.7);
\draw[thick] (0,0) circle (0.7);
\fill (0,0) circle (2.2pt);
\node[below left] at (0.1,0) {$O$};
\node at (-0.15,-1.25) {$(M)$};

% Point P
\fill (4.8,1.4) circle (2.2pt);
\node[above right] at (4.8,1.4) {$P$};
\node[right] at (5.0,1.0) {$(m)$};

% Segment OP
\draw[thick, dashed] (0,0) -- (4.8,1.4);
\node[below] at (2.45,0.55) {$R$};

% Vecteur unitaire u (de O vers P)
\draw[->, thick, blue] (0,0) -- (1.1,0.32);
\node[blue, above] at (0.9,0.45) {$\vec{u}$};

% Force vers O
\draw[->, very thick, red] (4.8,1.4) -- (3.35,0.98);
\node[red, above] at (4.0,1.35) {$\vec{F}$};

% Champ g vers O
\draw[->, very thick, green!50!black] (4.55,1.05) -- (3.6,0.78);
\node[green!50!black, below] at (4.05,0.82) {$\vec{g}(P)$};


\end{tikzpicture}
}%
\end{center}
}

% ============================
% ============================
% ============================
% Schéma : satellite en mouvement circulaire
% ============================
\newcommand{\SchemaSatelliteCirculaire}{%
\begin{center}
\resizebox{0.5\linewidth}{!}{%
\begin{tikzpicture}[>=Latex, every node/.style={font=\small}]

% ----------------------------
% Paramètres
% ----------------------------
\def\R{2.4}      % rayon de l'orbite
\def\ang{45}     % angle du point P sur l'orbite
\def\Ra{0.5}     % rayon de l'astre

% ----------------------------
% Centre et astre
% ----------------------------
\fill[orange!35] (0,0) circle (\Ra);
\draw[thick, orange!70!black] (0,0) circle (\Ra);

% Point O + label décalé
\fill (0,0) circle (2.2pt);
\node[below left] at (0.1,0) {$O$};

% ----------------------------
% Orbite circulaire
% ----------------------------
\draw[thick, dashed, gray!60] (0,0) circle (\R);

% ----------------------------
% Point satellite P sur le cercle
% ----------------------------
\coordinate (P) at (\ang:\R);
\fill (P) circle (2.2pt);
\node[above right] at (P) {$P$};

% ----------------------------
% Segment OP en pointillés
% ----------------------------
\draw[thick,dashed] (0,0) -- (P);
\node[below right] at ($(0,0)!0.58!(P)$) {$R$};

% ----------------------------
% Vecteur u radial
% ----------------------------
\draw[->, thick, blue] (0,0) -- ++({\ang}:0.9);
\node[blue, left, xshift=-1.5pt] at ({\ang}:0.8) {$\vec{u}$};

% ----------------------------
% Vitesse tangentielle
% ----------------------------
\draw[->, thick, green!50!black] (P) -- ++({\ang+90}:0.95);
\node[green!50!black, above left] at ($(P)+({\ang+90}:0.85)$) {$\vec{v}$};

% ----------------------------
% Vecteur u_t tangent
% ----------------------------
\draw[->, thick, cyan!60!blue] (P) -- ++({\ang+90}:0.70);
\node[cyan!60!blue, left] at ($(P)+({\ang+90}:0.2)$) {$\vec{u_t}$};

% ----------------------------
% Vecteur u_n vers le centre
% ----------------------------
\draw[->, thick, cyan!60!blue] (P) -- ++({\ang+180}:0.80);
\node[cyan!60!blue, right] at ($(P)+({\ang+180}:0.55)$) {$\vec{u_n}$};

\end{tikzpicture}%
}%
\end{center}%
}



% ============================
% Schéma : orbite elliptique (1re loi de Kepler)
% ============================
\newcommand{\SchemaEllipseKepler}{%
\begin{center}
\resizebox{0.5\linewidth}{!}{%
\begin{tikzpicture}[>=Latex, every node/.style={font=\small}]

\def\a{3.2}
\def\b{1.9}
\def\c{2.0}
\def\t{50}

% Ellipse
\draw[thick, gray!70] (0,0) ellipse [x radius=\a, y radius=\b];

% Grand axe
\draw[magenta!70,dashed] (-\a,0) -- (\a,0);
\node[below] at (0,-0.05) {grand axe};

% Centre
\fill (0,0) circle (1.8pt);
\node[below] at (0,0.5) {$O$};

% Foyers
\fill (-\c,0) circle (1.8pt);
\fill (\c,0) circle (1.8pt);
\node[above] at (-\c,0.2) {$F$};
\node[above] at (\c,0) {$F'$};

% Astre au foyer
\draw[orange!70!black] (-\c,0) circle (0.22);     % cercle creux
\fill[orange!70!black] (-\c,0) circle (0.05);     % petit point central
\node[below left] at (-\c,-0.08) {astre};

% Point P sur l'ellipse
\coordinate (P) at ({\a*cos(\t)},{\b*sin(\t)});
\fill (P) circle (2pt);
\node[above right] at (P) {$P$};

% Segments vers les foyers
\draw[dashed] (-\c,0) -- (P);
\draw[dashed] (\c,0) -- (P);

% Grand axe
\draw[<->, magenta!70, thick] (-\a,-2.15) -- (\a,-2.15);
\node[below, magenta!70] at (0,-2.15) {$2a$};

\end{tikzpicture}%
}%
\end{center}%
}

% ============================
% Schéma : 2e loi de Kepler (aires balayées)
% ============================
\newcommand{\SchemaAiresKepler}{%
\begin{center}
\resizebox{0.5\linewidth}{!}{%
\begin{tikzpicture}[>=Latex, every node/.style={font=\small}]

\def\a{3.3}
\def\b{2.0}
\def\c{2.1}

% Ellipse
\draw[thick, gray!70] (0,0) ellipse [x radius=\a, y radius=\b];

% Foyer
\fill (-\c,0) circle (2pt);
\node[below left] at (-\c,0) {$F$};

% Points sur l'ellipse
\coordinate (P1) at ({\a*cos(-130)},{\b*sin(-130)});
\coordinate (P2) at ({\a*cos(-35)},{\b*sin(-35)});
\coordinate (Q1) at ({\a*cos(160)},{\b*sin(160)});
\coordinate (Q2) at ({\a*cos(125)},{\b*sin(125)});

% Aires balayées
\fill[blue!20] (-\c,0) -- (P1) arc[start angle=-130,end angle=-35,x radius=\a,y radius=\b] -- cycle;
\fill[green!20] (-\c,0) -- (Q1) arc[start angle=160,end angle=125,x radius=\a,y radius=\b] -- cycle;

\draw[thick] (-\c,0) -- (P1);
\draw[thick] (-\c,0) -- (P2);
\draw[thick] (-\c,0) -- (Q1);
\draw[thick] (-\c,0) -- (Q2);

\fill (P1) circle (1.8pt);
\fill (P2) circle (1.8pt);
\fill (Q1) circle (1.8pt);
\fill (Q2) circle (1.8pt);

\node[blue!60!black] at (-0.2,-0.95) {$S_1$};
\node[green!50!black] at (-2.55,0.45) {$S_2$};

\node[right] at (P2) {$P_2$};
\node[right] at (P1) {$P_1$};
\node[left] at (Q1) {$P'_2$};
\node[left] at (Q2) {$P'_1$};

\end{tikzpicture}%
}%
\end{center}%
}
% ============================
% Schéma : 3e loi de Kepler
% ============================
\newcommand{\SchemaTroisiemeKepler}{%
\begin{center}
\resizebox{0.5\linewidth}{!}{%
\begin{tikzpicture}[>=Latex, every node/.style={font=\small}]

% Paramètres
\def\ap{3.4}   % demi-grand axe grande orbite
\def\bp{2.0}   % demi-petit axe grande orbite
\def\a{1.9}    % demi-grand axe petite orbite
\def\b{1.15}   % demi-petit axe petite orbite
\def\c{1.15}   % position du foyer commun

% Grande orbite
\draw[thick, magenta!80] (0,0) ellipse [x radius=\ap, y radius=\bp];

% Petite orbite
\draw[thick, olive!80!black] (\c,0) ellipse [x radius=\a, y radius=\b];

% Axe horizontal
\draw[gray!60] (-\ap,0) -- (\c+\a+0.2,0);

% Astre au foyer commun
\fill ( \c,0 ) circle (2pt);
\node[below right] at (\c,0) {$O$};

% Labels périodes
\node[magenta!80, font=\bfseries\large] at (-1.7,0.9) {Période $T'$};
\node[olive!80!black, font=\bfseries\large] at (1.15,0.35) {Période $T$};

% Demi-grand axe petite orbite : 2a
\draw[dashed, olive!80!black] (\c-\a,-0.02) -- (\c-\a,-1.45);
\draw[dashed, olive!80!black] (\c+\a,-0.02) -- (\c+\a,-1.45);
\draw[<->, thick, olive!80!black] (\c-\a,-1.20) -- (\c+\a,-1.20);
\node[olive!80!black, below] at (\c,-1.20) {$2a$};

% Demi-grand axe grande orbite : 2a'
\draw[dashed, magenta!80] (-\ap,-0.02) -- (-\ap,-2.25);
\draw[dashed, magenta!80] (\ap,-0.02) -- (\ap,-2.25);
\draw[<->, thick, magenta!80] (-\ap,-2.20) -- (\ap,-2.20);
\node[magenta!80, below] at (0,-2.20) {$2a'$};

\end{tikzpicture}%
}%
\end{center}%
}


% ============================
% ============================
% Schéma : orbite géostationnaire
% ============================
\newcommand{\SchemaGeostationnaire}{%
\begin{center}
\resizebox{0.65\linewidth}{!}{%
\begin{tikzpicture}[>=Latex, every node/.style={font=\small}]

\def\Rt{0.9}
\def\Rorb{4.0}
\def\tilt{-25}
\def\ang{-35}

% Orbite géostationnaire (vue en perspective)
\draw[thick, gray!60, rotate=\tilt] (0,0) ellipse [x radius=4.2, y radius=2.2];
\node[gray!70] at (-2.7,-1.45) {Orbite};
\node[gray!70] at (-2.7,-1.95) {géostationnaire};

% Flèche sur l'orbite
%\draw[->, gray!60, thick] (1.15,-2.15) arc[start angle=300,end angle=255,x radius=2.2,y radius=1.0];

% Terre
\fill[cyan!20] (0,0) circle (\Rt);
\draw[thick, cyan!50!black] (0,0) circle (\Rt);

% Centre de la Terre
\fill (0,0) circle (2pt);
\node[above left] at (-0.03,-0.07) {$T$};

% Axe de rotation incliné
\begin{scope}[rotate=-25]
  \draw[dashed, name path global=axe] (0,-1.5) -- (0,1.5);
\end{scope}

\node at (0.45,1.45) {Nord};
\node at (-0.3,-1.35) {Sud};


% Équateur en perspective
\begin{scope}[rotate=-25]
\draw[thick, red!70] (-\Rt,0) arc[start angle=180,end angle=360,x radius=\Rt,y radius=0.22];
\draw[thick, red!70, dashed] (\Rt,0) arc[start angle=0,end angle=180,x radius=\Rt,y radius=0.22];
\end{scope}
\node[red!70, scale=1.15] at (-1.75,0.35) {Équateur};

% Point sur l'orbite (sur l'ellipse tournée)
\begin{scope}[rotate=\tilt]
  \coordinate (S) at ({4.2*cos(\ang)},{2.2*sin(\ang)});
\end{scope}
\fill (S) circle (2pt);
\node[right] at (S) {$S$};

% Point sur la surface terrestre dans la direction de S
\coordinate (B) at ($(0,0)!\Rt/\Rorb!(S)$);

% Point S sur l'orbite
\begin{scope}[rotate=\tilt]
  \coordinate (S) at ({4.2*cos(\ang)},{2.2*sin(\ang)});
\end{scope}
\fill (S) circle (2pt);
\node[right] at (S) {$S$};

% Droite TS
\path[name path global=TSline] (0,0) -- ($(S)!1.2!(0,0)$);

% Intersection entre TS et l'axe de rotation
\path[name intersections={of=TSline and axe, by=I}];

% Point B sur la surface terrestre, dans la direction TS
\coordinate (B) at ($(0,0)!0.25!(S)$);

% RT : de l'axe de rotation à la surface
\draw[<->, thick, purple!80!black] (I) -- (B);
\node[purple!80!black, above right] at ($(I)!0.2!(B)$) {$R_T$};

% h_gs : de la surface à l'orbite
\draw[<->, thick, green!70!black] (B) -- (S);
\node[green!70!black, above] at ($(B)!0.58!(S)$) {$h_{gs}$};

% r_gs : segment parallèle décalé sous h_gs
\coordinate (Bbis) at ($(I)+(-0.18,-0.18)$);
\coordinate (Sbis) at ($(S)+(-0.18,-0.18)$);
\draw[<->, thick, green!70!black] (Bbis) -- (Sbis);
\node[green!70!black, below] at ($(Bbis)!0.5!(Sbis)$) {$r_{gs}$};


\end{tikzpicture}%
}%
\end{center}%
}
\begin{document}

\TitreBandeau{Mécanique céleste et satellites -- Fiche de cours}

\columnratio{0.5,0.5}
\setlength{\columnsep}{1.2cm}
\setlength{\columnseprule}{0.6pt}

\begin{paracol}{2}

% =====================================================
% ================= COLONNE GAUCHE ====================
% =====================================================

\section{Accélération d'un satellite}

\subsection{Champ de gravitation créé par un astre}

Soit un astre de masse $M$, à symétrie sphérique, assimilé à un point matériel placé en son centre de masse $O$.

Un point matériel $P$ de masse $m$, situé à la distance $R=OP$ du centre de l'astre, subit de la part de cet astre une force gravitationnelle attractive.

\prop La force gravitationnelle exercée par l'astre sur l'objet est :

\[
\vec{F}=-G\frac{Mm}{R^2}\vec{u}
\]

où :

\begin{itemize}
\item $G = 6{,}67\times10^{-11}\ \mathrm{N\,m^2\,kg^{-2}}$
\item $R$ est la distance entre les centres
\item $\vec{u}$ est un vecteur unitaire radial
\end{itemize}

\SchemaGravitation


\subsection{Champ gravitationnel}

Le champ gravitationnel créé par l'astre au point $P$ est noté $\vec{g}(P)$.


\switchcolumn
%\noindent
\vspace*{1\baselineskip}
\setcounter{section}{1}
\setcounter{subsection}{2}

\defi Le champ gravitationnel est relié à la force gravitationnelle par :

\[
\vec{F}=m\vec{g}(P)
\]


\prop L'expression du champ gravitationnel au point $P$ est donc :

\[
\boxed{\vec{g}(P)=-G\frac{M}{R^2}\vec{u}}
\]

\subsection{Accélération d'un satellite}

Un satellite est un système en mouvement autour d'un astre attracteur sous l'effet du champ gravitationnel créé par celui-ci.

Dans le référentiel d'étude, on applique la deuxième loi de Newton :

\[
m\vec{a}=\vec{F}
\]

or $\vec{F}=m\vec{g}(P)$, donc :

\[
\vec{a}=\vec{g}(P)
\]

\prop L'accélération du satellite est :

\[
\boxed{\vec{a}=-G\frac{M}{R^2}\vec{u}}
\]

Elle est radiale et centripète (dirigée vers le centre de l'astre).

% =====================================================
% ================= COLONNE DROITE ====================
% =====================================================

\switchcolumn
\setcounter{section}{1}

\section{Satellite en mouvement circulaire}

\subsection{Utilisation du repère de Frenet}

Soit un satellite en mouvement circulaire autour du centre $O$ de l'astre attracteur de masse $M$.

La distance $R=OP$ est constante.

On utilise le repère de Frenet :

\begin{itemize}
\item $\vec{u_t}$ : vecteur unitaire tangent à la trajectoire
\item $\vec{u_n}$ : vecteur unitaire normal dirigé vers le centre
\end{itemize}

\SchemaSatelliteCirculaire

En notant $v(t)$ la vitesse du satellite, l'accélération s'écrit :

\[
\vec{a}=\frac{dv}{dt}\vec{u_t}+\frac{v^2}{R}\vec{u_n}
\]

\subsection{Mouvement circulaire uniforme}

Pour un mouvement circulaire uniforme, la valeur de la vitesse est constante, donc :

\[
\frac{dv}{dt}=0
\]

\switchcolumn
\vspace*{1\baselineskip}
\setcounter{section}{2}
\setcounter{subsection}{2}

L'accélération devient alors :

\[
\vec{a}=\frac{v^2}{R}\vec{u_n}
\]

Or, d'après l'étude gravitationnelle précédente :

\[
\vec{a}=-G\frac{M}{R^2}\vec{u}
\]

Comme $\vec{u_n}$ et $-\vec{u}$ ont même sens, on identifie les normes :

\[
\frac{v^2}{R}=G\frac{M}{R^2}
\]

\subsection{Vitesse d'un satellite}

On en déduit :

\[
v^2=\frac{GM}{R}
\]

\prop La vitesse d'un satellite en orbite circulaire est :

\[
\boxed{v=\sqrt{\frac{GM}{R}}}
\]

\subsection{Période de révolution}

La période de révolution $T$ d'un satellite est la durée qu'il met pour effectuer un tour complet autour de l'astre attracteur.

Pour un mouvement circulaire de rayon $R$, la distance parcourue en un tour est  le périmètre du grand cercle :

\[
2\pi R
\]

Comme la vitesse est constante :

\[
v=\frac{2\pi R}{T}
\]

En remplaçant $v$ par $\sqrt{\dfrac{GM}{R}}$, on obtient :

\[
\frac{2\pi R}{T}=\sqrt{\frac{GM}{R}}
\]

\switchcolumn
\vspace*{0.4\baselineskip}
\setcounter{section}{2}
\setcounter{subsection}{4}


puis :

\[
T^2=\frac{4\pi^2r^3}{GM}
\]

\prop La période de révolution vérifie :

\[
\boxed{T=2\pi\sqrt{\frac{R^3}{GM}}}
\]

\rem Cette relation pourra être reliée ensuite aux lois de Kepler.


% =====================================================
% ================= COLONNE GAUCHE ====================
% =====================================================

\section{Première loi de Kepler}

\subsection{Loi des orbites}

\prop Les orbites des satellites sont des ellipses dont le centre de l'astre attracteur occupe un des foyers.

\SchemaEllipseKepler

\rem Une orbite circulaire est un cas particulier d'ellipse.

\subsection{Vocabulaire}

\begin{itemize}
\item Une ellipse possède 2 foyers, définie géométriquement comme tous les points P tel que FP +F'P=2a.
\item Le centre de l'astre attracteur se trouve en un foyer.
\item Le plus grand diamètre de l'ellipse est appelé \textbf{grand axe}.
\item Sa moitié est le \textbf{demi-grand axe}, noté $a$.
\end{itemize}

\switchcolumn
\setcounter{section}{3}


\section{Deuxième loi de Kepler}

\subsection{Loi des aires}

\prop Le segment reliant le centre de l'astre attracteur au satellite balaie des surfaces d'aires égales pendant des durées égales.

\SchemaAiresKepler

\rem Cela signifie que le satellite ne se déplace pas toujours à la même vitesse sur une orbite elliptique.

Il va plus vite lorsqu'il est proche de l'astre attracteur, et plus lentement lorsqu'il en est éloigné.

% =====================================================
% ================= p3 ====================
% =====================================================


\section{Troisième loi de Kepler}

\subsection{Loi des périodes}

\prop Le quotient du carré de la période T de révolution d'un satellite par le cube du demi-grand axe $a$ de son orbite est indépendant du satellite et ne dépend que de la masse $M$ de l'astre attracteur.

\[
\boxed{\frac{T^2}{a^3}={constante}}
\]

\SchemaTroisiemeKepler

\switchcolumn
\setcounter{section}{5}
\setcounter{subsection}{1}

\rem Pour tous les satellites en orbite autour d'un même astre attracteur, la quantité $\dfrac{T^2}{a^3}$ est la même.


\subsection{Cas particulier d'une orbite circulaire}

Si l'orbite est circulaire, alors le rayon $R$ de l'orbite joue le rôle du demi-grand axe :

\[
a=R
\]

On montre que la troisième loi de Kepler s'écrit aussi :

\[
\boxed{\frac{T^2}{R^3}=\frac{4\pi^2}{GM}}
\]


\section{Application : satellite géostationnaire}

\subsection{Définition}

\prop Un satellite géostationnaire est un satellite artificiel de la Terre fixe dans le référentiel terrestre.

\subsection{Conditions}

Pour qu'un satellite soit géostationnaire :

\begin{itemize}
\item son orbite doit être circulaire ;
\item son orbite doit être dans le plan de l'équateur ;
\item sa période doit être égale au jour sidéral :
\[
T = T_{\mathrm{sid}} = 86\,164\ \text{s}
\]
\end{itemize}

\switchcolumn
\setcounter{section}{6}
\setcounter{subsection}{2}

\SchemaGeostationnaire

\subsection{Application}

En appliquant la troisième loi de Kepler :

\[
\frac{T^2}{r_{gs}^3}=\frac{4\pi^2}{GM_T}
\]

on obtient pour l'orbite géostationnaire :

\[
r_{gs}\approx 4{,}2\times10^7\ \text{m}
\]

et donc :

\[
h_{gs}=r_{gs}-R_T \approx 3{,}6\times10^4\ \text{km}
\]

Un satellite géostationnaire est donc situé à environ \(36\,000\) km au-dessus de l'équateur.


\end{paracol}

\end{document}
